\section{Kiến Trúc Tổng Quát}
\label{sec:architecture}

Kiến trúc của \tool được tổ chức theo mô hình \textit{đồ thị trạng thái} (state graph), triển khai trên LangGraph --- một thư viện điều phối tác tử hỗ trợ đồ thị có hướng với các cạnh điều kiện (conditional edges), checkpoint trạng thái, và cơ chế truyền phát cập nhật theo thời gian thực. Pipeline bao gồm sáu nút (nodes) tương ứng với sáu tác tử chuyên biệt cùng một nút xử lý lỗi, được nhóm thành ba giai đoạn logic:

\begin{enumerate}
    \item \textbf{Giai đoạn Khởi tạo} (Initialization): Tác tử Giao tiếp ($\mathcal{A}_{conv}$) thu thập yêu cầu từ người dùng qua ngôn ngữ tự nhiên, trích xuất ý định dự đoán, và xác thực kết nối dữ liệu.
    \item \textbf{Giai đoạn Phân tích và Kỹ nghệ Dữ liệu} (Data Planning \& Engineering): Tác tử Phân tích ($\mathcal{A}_{eda}$) khám phá cấu trúc cơ sở dữ liệu và thực hiện phân tích dữ liệu khám phá; Tác tử Xây dựng Dữ liệu ($\mathcal{A}_{builder}$) chuyển đổi cơ sở dữ liệu quan hệ $(\mathcal{T}, \mathcal{L})$ sang đồ thị thực thể quan hệ (Mục~\ref{sec:relation-graph}); Tác tử Xây dựng Tác vụ ($\mathcal{A}_{task}$) sinh truy vấn SQL có nhúng cửa sổ thời gian để tạo bảng huấn luyện $T_{train}$ (Mục~\ref{sec:task_to_training_table}).
    \item \textbf{Giai đoạn Mô hình hóa và Thực thi} (Modeling \& Execution): Tác tử Chuyên gia GNN ($\mathcal{A}_{gnn}$) lựa chọn siêu tham số từ các nguồn tri thức ngoài và sinh kịch bản huấn luyện hiện thực hóa kiến trúc GNN thời gian (Mục~\ref{sec:encoders}); Tác tử Vận hành ($\mathcal{A}_{ops}$) thực thi huấn luyện trong môi trường cô lập và trích xuất kết quả đánh giá.
\end{enumerate}

\paragraph{Điều phối dựa trên đồ thị trạng thái.}
Trình điều phối (\texttt{PlexeOrchestrator}) xây dựng một \texttt{StateGraph} với sáu nút chức năng và một nút \texttt{error\_handler}. Các chuyển tiếp giữa các nút được kiểm soát bởi các \textit{hàm định tuyến có điều kiện} (conditional routing functions) --- ví dụ, hàm \texttt{route\_from\_conversation} kiểm tra sự tồn tại của nguồn dữ liệu và xác nhận người dùng trước khi chuyển sang giai đoạn phân tích; hàm \texttt{route\_from\_schema} xác minh rằng thư mục CSV chứa dữ liệu hợp lệ trước khi tiến đến bước xây dựng dữ liệu. Cơ chế này cho phép pipeline tự điều hướng linh hoạt: quay lại bước trước khi gặp lỗi, hoặc kết thúc sớm khi điều kiện tiên quyết không được đáp ứng, thay vì chạy tuần tự cứng nhắc. Trạng thái pipeline được lưu trữ qua \texttt{MemorySaver} checkpoint, cho phép tiếp tục phiên làm việc đã bị gián đoạn.

\paragraph{Bộ nhớ trạng thái toàn cục.}
Để đảm bảo tính nhất quán xuyên suốt quy trình, hệ thống duy trì một \textbf{bộ nhớ trạng thái toàn cục} $\mathcal{S}$ (biểu diễn qua \texttt{PipelineState} \texttt{TypedDict}), đóng vai trò là không gian lưu trữ chung giữa các tác tử. Bộ nhớ $\mathcal{S}$ bao gồm bốn nhóm thành phần:

\begin{itemize}
    \item \textbf{Lịch sử tương tác} $\mathcal{S}_{hist}$: Chuỗi các thông điệp ngôn ngữ tự nhiên giữa người dùng và hệ thống, được quản lý theo chiến lược \textit{append-only} (sử dụng \texttt{Annotated[List, operator.add]} trong LangGraph).
    \item \textbf{Tri thức nghiệp vụ} $\mathcal{S}_{schema}$: Đặc tả lược đồ cơ sở dữ liệu (\texttt{SchemaInfo}), bao gồm các ràng buộc khóa chính, khóa ngoại, nhãn thời gian, kết quả phân tích thống kê (\texttt{EDAInfo}), và bản phân loại bảng thực thể/sự kiện.
    \item \textbf{Đặc tả kỹ thuật} $\mathcal{S}_{spec}$: Ý định dự đoán của người dùng (\texttt{user\_intent}), thông tin tập dữ liệu (\texttt{DatasetInfo}) chứa dấu thời gian val/test, thông tin tác vụ (\texttt{TaskInfo}), và đường dẫn đến các tệp mã được sinh (\texttt{dataset.py}, \texttt{task.py}).
    \item \textbf{Trạng thái tiến trình} $\mathcal{S}_{exec}$: Nhật ký lỗi phân loại (\texttt{ErrorRecord}), bộ đếm thử lại có ngưỡng cấu hình riêng cho từng nút, pha hiện tại (\texttt{PipelinePhase}), và kết quả huấn luyện (\texttt{TrainingResult}).
\end{itemize}

Cơ chế bộ nhớ dùng chung theo kiến trúc \textit{blackboard}~\cite{hayes1985blackboard} loại bỏ nhu cầu giao tiếp điểm-đến-điểm (point-to-point) giữa các tác tử, từ đó giảm độ phức tạp tích hợp từ $O(n^2)$ xuống $O(n)$ khi số lượng tác tử tăng lên. Mỗi tác tử chỉ cần đọc các trường liên quan trong $\mathcal{S}$ và ghi lại kết quả đầu ra.

\paragraph{Mô hình tác tử.}
Mỗi tác tử kế thừa từ lớp trừu tượng \texttt{BaseAgent} cung cấp ba thành phần cốt lõi: (i)~một mô hình ngôn ngữ lớn (LLM) cấu hình riêng cho từng vai trò thông qua \texttt{AgentConfig}, (ii)~một tập hợp các công cụ (\textit{tools}) được đăng ký qua LangChain, và (iii)~một \textit{system prompt} nhúng tri thức chuyên gia dành riêng cho giai đoạn mà tác tử phụ trách. Kiến trúc này cho phép thay đổi mô hình nền tảng (ví dụ, GPT-4o, Claude, hoặc Gemini) cho từng tác tử riêng lẻ qua biến môi trường, mở ra khả năng tối ưu hóa chi phí-chất lượng theo vai trò. Ngoài ra, mỗi tác tử tự động tích hợp các công cụ từ các máy chủ MCP (Model Context Protocol) thông qua \texttt{MCPManager}, cho phép mở rộng bộ công cụ mà không cần sửa đổi mã nguồn tác tử.

Hình~\ref{fig:architecture} minh họa luồng điều khiển tổng thể của hệ thống.

\begin{figure}[htbp]
    \centering
    % \includegraphics[width=\linewidth]{figures/relau2ml_architecture.pdf}
    \caption{Kiến trúc tổng quan của \tool. Sáu tác tử chuyên biệt cùng một nút xử lý lỗi được tổ chức trong đồ thị trạng thái LangGraph. Mỗi tác tử đọc và ghi vào bộ nhớ trạng thái toàn cục $\mathcal{S}$; các mũi tên liền thể hiện luồng chính, các mũi tên đứt thể hiện luồng phục hồi lỗi.}
    \label{fig:architecture}
\end{figure}